\scanpage{198}

% B01 / text / top, powers of a generator
En termes d'un générateur $\unsure{\omega|w}_0$ de $\Lambda_E$,
i.e. $\omega_0\in\Lambda_E^*$, on trouve la structure de
contour orienté correspondant à $\unsure{\nu|v}\cdot[\omega_0]=[\omega_0^\nu]$,
via $u=\omega_0^\nu$ -- et pour $\dot\nu=\{\pm\nu\}$, on
prend la structure de contour définie par la
paire $\{\omega_0^\nu,\omega_0^{-\nu}\}$.

% B02 / text / upper middle, return to the situation
Revenons maintenant au cas de la situation
$(S,G_0,\subset\unsure{\mathfrak S|\mathfrak G}_S,\ \unsure{\mathcal E|\mathscr E|\Sigma}\subset\unsure{P|\mathcal P|\mathfrak P}_E)$, et on voit donc que les
systèmes
% slanted
\margin{\uline{NB} On voit maintenant que $G$ est transitif sur l'ens des}

% B03 / display / middle, system of contours
\[
 (C_E)_{E\in\mathcal E}\qquad
 \begin{aligned}
 &C_E\in\text{ structures de}\\
 &\hspace{2em}\text{contour sur }E\\
 &C_E=(E,A_E)
 \end{aligned}
\]

% B04 / text / middle, parametrization, left margin continuation
invariés par $G$, correspondent biunivoquement
aux éléments \inserted{\deleted{\unsure{$u$}}$\unsure{v|\nu}_{E_0}$} de
$\bigl(\Lambda_{E_0}\smallsetminus\underbrace{\{0,\sigma\}}_{\text{sous-groupe d'ordre 2}}\bigr)/\mathord{\pm}1$.
% slanted
\margin{ou \unsure{au|aux|on} syst. d'éléments $(v_E)$ \unsure{aux|avec} $v_E\in(\Lambda_E\smallsetminus\{0,\sigma_E\})/\mathord{\pm}1$ \unsure{associés|invariés} par $G\ldots$}
\deleted{\unsure{En}}\deleted{Nous} \unsure{$\ldots$}

% B05 / text / lower middle, realization question
Nous proposons maintenant, \deleted{pour} un tel
$v_{E_0}$ ou $v_E$, d'examiner s'il provient bien
d'une \inserted{structure} $i\in\Pi$.

% B06 / text / lower, assigning the edge set
\deleted{Supposons donc}
Considérons donc le système \deleted{$(C_E)$}
$C_E=(E,A_E)$ associé à un $v_{E_0}$ (ou à un $(v_E)_{E\in\mathcal E}$)
et posons

% B07 / display / lower, (24)
\[
 A=\bigcup_{E\in\mathcal E} A_E\qquad\subset P_2(\unsure{S|B}), \tag*{(24)}
\]

% B08 / text / foot, transitivity
il est clair que $A$ est stable sous $G$,
et que $G$ est transitif sur l'ens des
arêtes orientées de la structure de graphe $(S,A)$.

% archive sheet number: 197
